% =============================================================
%  § 2. Шифрование: от античности до наших дней
% =============================================================

\section{Шифрование: от тайнописи к математической задаче}
\label{sec:history}

\subsection{Зачем людям шифры}

Каждый из вас, заходя на любой сайт по адресу, начинающемуся с \texttt{https://}, доверяет ему \emph{шифрование}. Браузер и сервер обмениваются миллионами сообщений, но никто посторонний прочесть их не может. В этом параграфе мы поймём, какие задачи решает шифрование, какие общие схемы существуют и почему именно теория чисел оказалась его математической основой.

\paragraph*{Античность.} Первые шифры были чисто механическими. В Древней Греции (V век до н.\,э.) воины Спарты использовали \term{скиталу} --- цилиндр определённого диаметра, на который наматывалась лента с текстом. Без точно такой же палочки прочесть сообщение было нельзя. Юлий Цезарь, по свидетельству Светония, в переписке со своими полководцами сдвигал каждую букву алфавита на три позиции: \texttt{A → D}, \texttt{B → E}, и так далее. Этот \term{шифр Цезаря} --- простейший \emph{шифр замены} (рис.~\ref{fig:p2-1}).

\begin{figure}[H]
\centering
\begin{tikzpicture}[scale=0.85, every node/.style={font=\small\ttfamily}]
  % Открытый текст
  \foreach \i/\ch in {0/V,1/E,2/N,3/I,4/V,5/I,6/D,7/I,8/V,9/I,10/C,11/I}{
    \draw[prosvBlue] (\i*0.6, 1.2) rectangle (\i*0.6+0.55, 1.7);
    \node at (\i*0.6+0.275, 1.45) {\ch};
  }
  \node[left, font=\small\rmfamily] at (0,1.45) {Открытый текст:};
  
  % Стрелки
  \foreach \i in {0,...,11}{
    \draw[->,>=stealth, prosvGray] (\i*0.6+0.275, 1.15) -- (\i*0.6+0.275, 0.75);
  }
  
  % Шифр-текст (сдвиг +3)
  \foreach \i/\ch in {0/Y,1/H,2/Q,3/L,4/Y,5/L,6/G,7/L,8/Y,9/L,10/F,11/L}{
    \draw[prosvRed] (\i*0.6, 0.2) rectangle (\i*0.6+0.55, 0.7);
    \node at (\i*0.6+0.275, 0.45) {\ch};
  }
  \node[left, font=\small\rmfamily] at (0,0.45) {Шифр-текст:};
  
  \node[right, font=\small\rmfamily] at (7.5,1.0) {\itshape\color{prosvGray}«пришёл, увидел, победил»};
\end{tikzpicture}
\caption{Шифр Цезаря со сдвигом~3: каждая буква заменяется на букву через три позиции в алфавите.}
\label{fig:p2-1}
\end{figure}

Любой такой шифр легко взламывается: букв в алфавите немного, и можно просто перебрать все 25 возможных сдвигов. Более изобретательные шифры (например, \term{шифр Виженера} XVI века) использовали \emph{переменный} сдвиг, заданный ключевым словом. Они продержались несколько столетий, пока в середине XIX века Чарлз Бэббидж и независимо Фридрих Касиски не предложили универсальный метод их взлома.

\paragraph*{XX век: машины.} Во время Второй мировой войны на сцену вышли электромеханические шифровальные машины --- знаменитая немецкая \term{Энигма}, японская «Пурпурная» (Purple), советская «Фиалка». Их взлом стал отдельной героической страницей: в Великобритании Алан Тьюринг возглавлял проект \emph{Bletchley Park}, где была сконструирована машина «Бомба», взламывавшая Энигму. Считается, что этот успех сократил войну в Европе на 2--4 года.

\begin{interestbox}
Деятельность Алана Тьюринга и его коллег была настолько секретной, что заслуги команды Bletchley Park были рассекречены только в 1970-х годах --- то есть через 25 лет после конца войны. Сам Тьюринг при жизни так и не узнал о признании. Тьюринг также считается одним из отцов теоретической информатики --- именно его именем названа престижная Премия Тьюринга, аналог «Нобелевской премии в информатике».
\end{interestbox}

% -----------------------------------------------------------------
\subsection{Общая схема симметричного шифрования}

Все шифры до 1970-х годов имели одну общую особенность: для шифрования и расшифрования использовался \emph{один и тот же} ключ. Такие шифры называются \term{симметричными} (или \engterm{secret-key}) (рис.~\ref{fig:symenc}).

\begin{figure}[H]
\centering
\begin{tikzpicture}[
  font=\small,
  every node/.style={align=center},
  box/.style={rectangle, draw=prosvBlue, fill=prosvLight, minimum width=2.5cm, minimum height=1cm, rounded corners=2pt},
  evil/.style={rectangle, draw=prosvRed, fill=prosvCream, minimum width=2.2cm, minimum height=0.8cm, rounded corners=2pt}
]
  \node[circle, draw=prosvBlue, fill=white, minimum size=0.7cm] (a) at (-5,0) {A};
  \node[below, font=\scriptsize] at (-5,-0.5) {Алиса};
  \node[circle, draw=prosvBlue, fill=white, minimum size=0.7cm] (b) at (5,0) {B};
  \node[below, font=\scriptsize] at (5,-0.5) {Боб};
  
  \node[box] (enc) at (-2.5,0) {Шифрование $E_K$};
  \node[box] (dec) at (2.5,0) {Расшифрование $D_K$};
  
  \draw[->,>=stealth, prosvBlue, thick] (a) -- node[above]{$m$} (enc);
  \draw[->,>=stealth, prosvBlue, thick] (enc) -- node[above]{$c = E_K(m)$} (dec);
  \draw[->,>=stealth, prosvBlue, thick] (dec) -- node[above]{$m$} (b);
  
  % Ключи
  \node[font=\Large] at (-2.5, 1.2) {[K]};
  \node[font=\scriptsize] at (-2.5, 1.7) {ключ $K$};
  \node[font=\Large] at (2.5, 1.2) {[K]};
  \node[font=\scriptsize] at (2.5, 1.7) {ключ $K$};
  \draw[dashed, prosvGray] (-2.5,1.0) -- (2.5,1.0);
  
  % Ева
  \node[evil] (eve) at (0, -1.8) {Ева подслушивает\\ канал};
  \draw[->,>=stealth, prosvRed, dashed] (0, -0.3) -- (eve);
\end{tikzpicture}
\caption{Симметричное шифрование: один общий секретный ключ $K$ позволяет и зашифровать, и расшифровать. Подслушивающая Ева не должна узнать $K$.}
\label{fig:symenc}
\end{figure}

Симметричные шифры быстры и до сих пор широко применяются (например, стандарт \engterm{AES} --- то, чем шифруется содержимое любого Telegram-сообщения после установки соединения). Но у них есть фундаментальная \emph{проблема распределения ключей}: как Алисе и Бобу договориться о ключе, если они никогда раньше не встречались, а связь между ними прослушивается?

\begin{importantbox}
\textbf{Парадокс симметричного шифрования.} Чтобы безопасно обмениваться сообщениями, нужен общий ключ. Чтобы безопасно передать ключ --- нужен другой общий ключ. И так до бесконечности.
\end{importantbox}

В большом мире (миллиарды пользователей интернета) задача распределения ключей становится непреодолимой: каждой паре нужен свой ключ, всего $\binom{n}{2} \approx n^2/2$ ключей.

% -----------------------------------------------------------------
\subsection{Революция 1976 года: открытый ключ}

В 1976 году произошёл прорыв, перевернувший всю криптографию. Уитфилд Диффи и Мартин Хеллман предложили принципиально новую идею: \term{асимметричное шифрование}, или \term{криптография с открытым ключом}.

Каждый пользователь имеет \emph{пару} ключей:
\begin{itemize}
  \item \term{открытый ключ} (\engterm{public key}) --- его можно (и нужно) сообщать всем: публиковать на сайте, отправлять незашифрованно;
  \item \term{закрытый ключ} (\engterm{private key}) --- он хранится в строжайшем секрете у владельца.
\end{itemize}

Ключевое свойство: что зашифровано открытым ключом --- расшифровать может только владелец соответствующего закрытого ключа (рис.~\ref{fig:asymenc}).

\begin{figure}[H]
\centering
\begin{tikzpicture}[
  font=\small,
  every node/.style={align=center},
  box/.style={rectangle, draw=prosvBlue, fill=prosvLight, minimum width=2.4cm, minimum height=1cm, rounded corners=2pt}
]
  \node[circle, draw=prosvBlue, fill=white, minimum size=0.7cm] (a) at (-5,0) {A};
  \node[below, font=\scriptsize] at (-5,-0.5) {Алиса};
  \node[circle, draw=prosvBlue, fill=white, minimum size=0.7cm] (b) at (5,0) {B};
  \node[below, font=\scriptsize] at (5,-0.5) {Боб};
  
  \node[box] (enc) at (-2.3,0) {Шифрование $E$};
  \node[box] (dec) at (2.3,0) {Расшифрование $D$};
  
  \draw[->,>=stealth, prosvBlue, thick] (a) -- node[above]{$m$} (enc);
  \draw[->,>=stealth, prosvBlue, thick] (enc) -- node[above]{$c$} (dec);
  \draw[->,>=stealth, prosvBlue, thick] (dec) -- node[above]{$m$} (b);
  
  % Открытый ключ Боба
  \node[font=\Large] at (-2.3, 1.2) {[\textcolor{prosvGreen}{pk}]};
  \node[font=\scriptsize, text=prosvGreen, align=center] at (-2.3, 1.85) {откр. ключ\\ Боба $\mathit{pk}_B$};
  
  % Закрытый ключ Боба
  \node[font=\Large] at (2.3, 1.2) {[\textcolor{prosvRed}{sk}]};
  \node[font=\scriptsize, text=prosvRed, align=center] at (2.3, 1.85) {закр. ключ\\ Боба $\mathit{sk}_B$};
  
  % Стрелка от Боба к открытому ключу — "раздаёт всем"
  \draw[->,>=stealth, prosvGreen, dashed] (b) to[bend right=40] (-2.3, 1.4);
  \node[text=prosvGreen, font=\scriptsize] at (1.5, 2.5) {публикуется};
\end{tikzpicture}
\caption{Асимметричное шифрование: Алиса использует \emph{открытый} ключ Боба, который тот заранее опубликовал. Расшифровать может только Боб --- его \emph{закрытым} ключом.}
\label{fig:asymenc}
\end{figure}

Заметьте: больше нет проблемы распределения ключей. Алиса не должна заранее ничего секретного передавать Бобу --- она просто берёт его открытый ключ из открытого источника.

% -----------------------------------------------------------------
\subsection{Односторонние функции и функции с секретом}

В чём математическая суть асимметричной криптографии? В существовании \term{односторонних функций}.

\begin{definition}{}{def:p2-1}
\term{Односторонняя функция} (\engterm{one-way function}) --- это функция $f \colon X \to Y$, обладающая двумя свойствами:
\begin{itemize}
  \item \emph{Прямое вычисление просто:} зная $x$, легко (полиномиально быстро) вычислить $f(x)$.
  \item \emph{Обращение трудно:} зная $y = f(x)$, восстановить $x$ без какой-либо дополнительной информации --- задача \emph{практически} нерешаемая (требует экспоненциального времени).
\end{itemize}
\end{definition}

\begin{interestbox}
Существование «настоящих» односторонних функций математически не доказано (это связано с одной из главных нерешённых задач --- проблемой $\mathbf{P}\ne\mathbf{NP}$). На практике мы используем функции, для которых обращение \emph{считается} трудным, потому что человечество за десятилетия так и не научилось делать это быстро.
\end{interestbox}

Хороший бытовой образ: \emph{разбить вазу} --- легко, \emph{собрать обратно} --- очень тяжело. Или: \emph{перемешать колоду карт} --- мгновенная операция, \emph{восстановить исходный порядок} --- немыслимо.

Главные примеры из теории чисел:
\begin{itemize}
  \item \term{Умножение vs. факторизация:} перемножить $p\cdot q$ --- мгновенно; разложить произведение на простые --- трудно. Это база RSA.
  \item \term{Дискретный логарифм:} возвести $g$ в степень $x$ по модулю $p$ --- быстро; зная $g^x \bmod p$, восстановить $x$ --- очень трудно. Это база схемы Диффи--Хеллмана.
\end{itemize}

\paragraph*{Функции с секретом.} Просто односторонней функции для асимметричного шифрования мало: нужно, чтобы у \emph{кого-то одного} (владельца) была возможность всё-таки обратить функцию --- именно для расшифрования сообщений (рис.~\ref{fig:p2-4}).

\begin{definition}{}{def:p2-2}
\term{Функция с секретом} (\engterm{trapdoor function}) --- односторонняя функция $f$, для которой существует дополнительная секретная информация $t$ (\emph{лазейка}, \emph{trapdoor}), знание которой делает обращение $f$ полиномиально быстрым.
\end{definition}

\begin{figure}[H]
\centering
\begin{tikzpicture}[font=\small, every node/.style={align=center}]
  % Дверь без ключа
  \node[draw=prosvBlue, fill=prosvLight, rounded corners=2pt, minimum width=4cm, minimum height=1cm] (f) at (0,0) {Прямое вычисление:\\ \textit{очень быстро}};
  \node[draw=prosvRed, fill=prosvCream, rounded corners=2pt, minimum width=4cm, minimum height=1cm] (fi) at (0,-1.8) {Обращение без ключа:\\ \textit{экспоненциально долго}};
  \node[draw=prosvGreen, fill=prosvLight, rounded corners=2pt, minimum width=4cm, minimum height=1cm] (fis) at (0,-3.6) {Обращение со знанием секрета $t$:\\ \textit{полиномиально быстро}};
  
  \draw[->,>=stealth, thick, prosvBlue] (-3.0,0) -- (-2.05,0);
  \node[left, font=\scriptsize, text=prosvGray] at (-3,0) {$x$};
  \draw[->,>=stealth, thick, prosvBlue] (2.05,0) -- (3,0);
  \node[right, font=\scriptsize, text=prosvGray] at (3,0) {$f(x) = y$};
  
  \draw[->,>=stealth, thick, prosvRed] (3,-1.8) -- (2.05,-1.8);
  \node[right, font=\scriptsize, text=prosvGray] at (3,-1.8) {$y$};
  \draw[->,>=stealth, thick, prosvRed] (-2.05,-1.8) -- (-3,-1.8);
  \node[left, font=\scriptsize, text=prosvGray] at (-3,-1.8) {$x$};
  
  \draw[->,>=stealth, thick, prosvGreen] (3,-3.6) -- (2.05,-3.6);
  \node[right, font=\scriptsize, text=prosvGray] at (3,-3.6) {$y + t$};
  \draw[->,>=stealth, thick, prosvGreen] (-2.05,-3.6) -- (-3,-3.6);
  \node[left, font=\scriptsize, text=prosvGray] at (-3,-3.6) {$x$};
\end{tikzpicture}
\caption{Функция с секретом ведёт себя как односторонняя, но «открывается» при знании дополнительной информации~$t$.}
\label{fig:p2-4}
\end{figure}

Можно представить себе сейф с замком: запереть сейф (положить документ внутрь) может каждый, но открыть --- только владелец ключа.

\paragraph*{Идея использования.} Открытый ключ \emph{задаёт} функцию с секретом $f$. Кто угодно может вычислить $c = f(m)$ --- зашифровать сообщение. Для расшифрования нужно обратить $f$, что без секрета $t$ --- безнадёжно. У владельца секрета $t$ это вычисление быстрое. Тем самым роль закрытого ключа --- хранить лазейку $t$.

% -----------------------------------------------------------------
\subsection{Где криптография работает прямо сейчас}

Чтобы дальнейший материал не казался академическим упражнением, перечислим, в каких знакомых каждому ситуациях работают (часто прямо сейчас, пока вы это читаете) описанные в этой главе идеи.

\begin{itemize}
  \item \textbf{Загрузка любого сайта} (\texttt{https}). Браузер договаривается с сервером о ключе через схему Диффи--Хеллмана, после чего весь трафик шифруется симметричным алгоритмом (\engterm{AES}).
  \item \textbf{Мессенджеры} (WhatsApp, Signal, Telegram). WhatsApp и Signal по умолчанию используют \emph{сквозное шифрование} (\engterm{end-to-end}): сервер видит лишь шифр-текст. В Telegram сквозное шифрование включается только в «секретных чатах», а обычные облачные чаты шифруются лишь между клиентом и сервером (и хранятся на серверах Telegram). Согласование ключей --- асимметричное.
  \item \textbf{Банковские карты и оплата.} Транзакция от карты проходит через цепочку обмена цифровыми подписями.
  \item \textbf{Госуслуги и налоговая.} Подача декларации сопровождается \term{электронной цифровой подписью} (ЭЦП), о которой подробно --- в §~\ref{sec:rsa-apps}.
  \item \textbf{Цифровой пропуск/QR-код} (от билета в кино до подтверждения вакцинации). Содержит подпись эмитента, проверяемую его открытым ключом.
  \item \textbf{Криптовалюты.} Биткойн и его «родственники» --- огромная распределённая система с электронными подписями: кошелёк --- это просто пара (открытый ключ, закрытый ключ).
\end{itemize}

\paragraph*{Что должно быть у «правильной» криптосистемы.} Когда мы строим конкретную схему, надо проверять, что она удовлетворяет следующим естественным требованиям:
\begin{enumerate}
  \item \emph{Корректность.} Если Алиса честно зашифровала открытым ключом Боба, Боб всегда расшифрует --- ровно исходное сообщение.
  \item \emph{Стойкость.} Подслушивающий, видя только шифр-текст и открытый ключ, не может практически восстановить исходное сообщение.
  \item \emph{Эффективность.} Шифрование и расшифрование выполняются быстро (полиномиально).
  \item \emph{Защита ключа.} Закрытый ключ не должен «утекать» из открытого --- даже зная открытый ключ, нельзя восстановить закрытый.
\end{enumerate}

В следующих параграфах мы построим конкретные схемы --- RSA и Диффи--Хеллман --- и убедимся, что они этим требованиям удовлетворяют.

\begin{tasksbox}
\begin{enumerate}
  \item Зашифруйте шифром Цезаря со сдвигом 5 слово \texttt{INFORMATIKA}.
  \item Почему симметричное шифрование плохо подходит для миллиарда пользователей? Подсчитайте, сколько ключей надо хранить.
  \item Объясните своими словами, в чём состоит свойство «односторонности» функции.
  \item* Знаменитая Энигма имела примерно $10^{20}$ возможных стартовых установок. Хватило бы перебора, если перебирать миллион вариантов в секунду?
  \item Приведите три бытовые операции, которые ведут себя как «односторонние функции».
\end{enumerate}
\end{tasksbox}
